%%%%%%%%%%  scan idx 16  |  folio 3  %%%%%%%%%%
\subsubsection*{2.1}
Pour ce qui est de cette dernière, une définition s'offre de façon assez évidente. On utilise le fait que le morphisme $f$, grâce à l'hypothèse d'existence d'un Module inversible ample sur $X$ et sur $Y$, peut se factoriser en
\begin{equation}\tag{2.1}
X\xrightarrow{i}X'\xrightarrow{f'}Y,
\end{equation}
où $i$ est une immersion fermée, et $f'$ un morphisme lisse~; par exemple on pourra prendre $X'=\mathbb P^r_Y$ pour $r$ convenable. On posera alors
\begin{equation}\tag{2.2}
\check T_f=\cl({\underline\Omega}^1_{X'/Y})-\cl({\underline N}_{X/X'}),
\end{equation}
où ${\underline\Omega}^1_{X'/Y}$ est le Module localement libre des $1$-différentielles relatives de $X'$ sur $Y$ (ou Module cotangent relatif de $X'$ sur $Y$), et ${\underline N}_{X/X'}$ le Module conormal de $X$ dans $X'$, égal par définition à ${\underline J}/{\underline J}^2$, où ${\underline J}$ est l'Idéal sur $X'$ définissant le sous-schémas fermé $X$~; comme $X$ et $X'$ sont réguliers, on en conclut d'ailleurs que ${\underline N}_{X/X'}$ est également localement libre. Enfin $\check T_f$ désigne le dual de $T_f$. On vérifie sans peine (Exp.\ VIII) que l'élément $T_f$ défini par (2.1) ne dépend pas de la factorisation de $f$ choisie.

\subsubsection*{2.2.}
Quant à une définition, sous les conditions envisagées, d'un anneau de Chow et d'une théorie des classes de Chern correspondante, bien qu'il soit plausible qu'une telle théorie doive pouvoir se développer sur le modèle de celle déjà connue pour les schémas algébriques lisses et quasi-porjectifs, il semble que ce travail n'ait pas été fait à l'heure actuelle, et à fortiori il n'a pas été entrepris dans le Séminaire. On introduit par contre un autre anneau, qui jouera un rôle analogue à celui de l'anneau de Chow, et dont le produit tensoriel par $Q$ est effectivement isomorphe à $A(X)\otimes_{\mathbb Z}\mathbb Q$ dans le cas où $X$ est quasi-projectif et lisse sur un corps $k$. Cet anneau est l'anneau gradué associé à $K(X)$, pour une filtration convenable de cet anneau. Une première filtration naturelle de $K(X)$, qui est celle envisagée dans [RRR], est la «~\underline{filtration topologique}~» par la codimension du support des faisceaux cohérents, $\operatorname{Filt}^{i}_{\mathrm{top}}K(X)$ étant engendré par les $\cl(F)$ avec $F$ cohérent et $\operatorname{codim}(\operatorname{supp}F,X)\ge i$. Une première question est celle de la compatibilité de cette filtration avec la structure d'anneau, question résolue dans [RRR], dans le cas $X$ lisse sur un coprs $k$, grâce au «~moving lemma~» de Chow. Ce point admis, on trouve donc un anneau gradué associé, que nous noterons ici $\Gr^{\bullet}_{\mathrm{top}}(X)$. Dans [RRR] on définit un homomorphisme d'anneaux surjectif
\[
\phi:A(X)\longrightarrow \Gr^{\bullet}_{\mathrm{top}}(X),
\]
